\documentclass[10pt,a4paper]{amsart}
\usepackage[margin=19mm]{geometry}
\usepackage{amsmath,amssymb,mathtools}
\usepackage[scr=rsfs,cal=euler]{mathalfa}
\usepackage{xfrac}
\usepackage[unicode,hidelinks,bookmarks=false]{hyperref}
\newcommand{\ie}{i.e.\ }
\newcommand{\cf}{cf.\ }
\newcommand{\resp}{respectively\ }
\newcommand{\Subsection}{\subsection}
\providecommand{\nobreakdash}{\nobreak\hbox{-}}
\setlength{\emergencystretch}{2em}
\multlinegap0pt
\usepackage{amssymb}
\usepackage{enumerate}
\usepackage{xcolor}
\usepackage{float}
\usepackage{mathtools}
\usepackage{tikz}
\newtheorem{lemma}{Lemma}
\newtheorem{theorem}{Theorem}
\usepackage{cleveref}
\crefname{lemma}{Lemma}{Lemmas}
\crefname{theorem}{Theorem}{Theorems}
\def\bbv{\mathbb{V}}
\def\bdyn{\mathbf{Dyn}}
\DeclareMathOperator{\supp}{supp}
\title{Discrimination of dynamic data via curvature sets}
\author{Nadezhda Belova, Maxwell Goldberg, Facundo Memoli, Sriram Raghunath, Andrew Xie}
\date{Source: arXiv:2603.04624v1 (CC BY 4.0)}
\begin{document}
\maketitle

\noindent\emph{Source and licence.} Discrimination of dynamic data via curvature sets. Version arXiv:2603.04624v1 (CC BY 4.0), distributed by arXiv under the Creative Commons Attribution 4.0 International licence, http://creativecommons.org/licenses/by/4.0/, as recorded in the arXiv licence metadata for this exact version and shown on the abstract page https://arxiv.org/abs/2603.04624v1. Full licence notice: \texttt{licenses/arxiv-2603.04624-CC-BY-4.0.txt}.\par\smallskip

\noindent\textit{Editorial scope.} Counted unit: the article's theorem thm:intervaldecomp (source0.tex lines 348--350). The article's complete original proof of it is delivered with it (source0.tex lines 351--372, one \texttt{proof} environment of 22 lines). No mathematics is written, repaired or re-derived: the statement and the proof are the article's own lines, in the article's own notation, and no retained line is substituted or re-flowed.  Retained uncounted as the source-local context the proof needs: lines 330--332.  Nothing was omitted from any retained span, so the package carries no omission record.  No retained line contains a citation, so the wrapper carries no bibliography and no bibliography entry.  No retained line contains a figure environment, an image-inclusion command or drawing code, so no image is delivered and the package declares no asset dependency.  Editorial formatting only, in addition: an AMS \texttt{amsart} preamble that restates the article's own declarations and macros used by the retained lines, namely the environments \texttt{lemma}, \texttt{theorem} on separate counters, together with the standard packages those lines need.  The retained environments print in this document as Lemma 1, Theorem 1 in order of appearance, whereas the article prints the same items with the numbering of its own full document.  The complete native source of the article as distributed in the arXiv e-print for this exact version is bundled unchanged as \texttt{sources/195750/source0.tex}.
    \begin{lemma}\label{lem:structuremaps}
        Suppose $k \in \mathbb{Z}_{\geq 0}$ and let $\gamma_X = (X, d_X(\cdot))$ be a $(2k+2)$-DMS. Put $\bbv = H_k(\mathcal{R}^{\mathrm{lev}}(\gamma_X))$. Then for any $p, q \in \supp \bbv$ satisfying $p \leq q$, we have $\mathcal{R}^{\mathrm{lev}}(\gamma_X)(p) = \mathcal{R}^{\mathrm{lev}}(\gamma_X)(q)$ as abstract simplicial complexes. Moreover, the structure map $\bbv(p \leq q)$ is an isomorphism.
    \end{lemma}

    \begin{theorem}[Interval Decomposability] \label{thm:intervaldecomp}
        Suppose $k \in \mathbb{Z}_{\geq 0}$, and $\gamma_X = (X, d_X(\cdot))$ is a $(2k+2)$-DMS. Then $H_k(\mathcal{R}^{\mathrm{lev}}(\gamma_X))$ is interval-decomposable. Moreover, $H_k(\mathcal{R}^{\mathrm{lev}}(\gamma_X)) \cong \bigoplus_{\alpha \in A} \mathbb{I}_{S_\alpha}$, where $\{S_\alpha\}_{\alpha \in A}$ are the order-connected components of $\supp H_k(\mathcal{R}^{\mathrm{lev}}(\gamma_X))$.
    \end{theorem}

    \begin{proof}
        Let $\bbv = H_k(\mathcal{R}^{\mathrm{lev}}(\gamma_X))$, and let $\{S_\alpha\}_{\alpha \in A}$ be the decomposition of $\supp \bbv$ into order-connected components. We will first show that $\bbv \cong \mathbb{I}_{(\supp \bbv)}$.

        By \Cref{lem:structuremaps}, for any $p, q \in \supp \bbv$ such that $p \leq q$, it follows that $\mathcal{R}^{\mathrm{lev}}(\gamma_X)(p) = \mathcal{R}^{\mathrm{lev}}(\gamma_X)(q)$ as simplicial complexes. Thus, for any fixed $\alpha \in A$, $\mathcal{R}^{\mathrm{lev}}(\gamma_X)(p) = \mathcal{R}^{\mathrm{lev}}(\gamma_X)(q)$ for all $p, q \in S_\alpha$. As the underlying simplicial complexes are equal, we can identify their chain complexes. For each $\alpha$, choose an arbitrary $p_\alpha \in S_\alpha$. Because $\bbv(p_\alpha) = H_k(\mathcal{R}^{\mathrm{lev}}(\gamma_X)(p_\alpha)) \cong \mathbb{F}$, we can pick a representative cycle $[\xi_\alpha]$ that generates $\bbv(p_\alpha)$. And thus $[\xi_\alpha]$ generates $\bbv(q)$, for all $q \in S_\alpha$.

        For any $p \in \bdyn$, we can define the maps $f_p: \bbv(p) \rightarrow \mathbb{I}_{(\supp \bbv)}(p)$ and $g_p: \mathbb{I}_{(\supp \bbv)}(p) \rightarrow \bbv(p)$ by
        \[
            f_p(x) = 
            \begin{cases}
                t & \exists!\, \alpha \in A \text{ s.t. } p \in S_\alpha, x = t [\xi_\alpha] \\
                0 & \text{otherwise}
            \end{cases}
            \qquad
            g_p(t) = 
            \begin{cases}
                t[\xi_\alpha] & \exists! \, \alpha \in A \text{ s.t. } p \in S_\alpha \\
                0 & \text{otherwise}
            \end{cases}
        \]

        Which satisfy naturality and are inverses of each other. Thus, $\bbv \cong \mathbb{I}_{(\supp \bbv)}$. From this, it follows that $\bbv \cong \mathbb{I}_{(\supp \bbv)} \cong \bigoplus_{\alpha \in A} \mathbb{I}_{S_\alpha}$, which is the interval decomposition of $\bbv$.
    \end{proof}

\end{document}
